
\section*{E21/22 --- Cadena celular de vacancias HMT}

Definimos
\[
\rho=\log_{1000}(729)=\log_{10}9,
\qquad
\epsilon=1-\rho=\log_{10}\frac{10}{9}.
\]
Entonces
\[
\rho=0.95424250943932487459\ldots,
\qquad
\epsilon=0.04575749056067512540\ldots.
\]
La palabra de vacancia es
\[
z_t=\lceil(t+1)\epsilon\rceil-\lceil t\epsilon\rceil.
\]
Para cualquier ventana:
\[
Q_N(m)=\sum_{t=m}^{m+N-1}z_t
=\lceil(m+N)\epsilon\rceil-\lceil m\epsilon\rceil,
\]
de modo que
\[
\boxed{Q_N(m)\in\{\lfloor N\epsilon\rfloor,\lceil N\epsilon\rceil\}.}
\]

Las posiciones de vacancia
\[
p_n=\left\lfloor\frac{n}{\epsilon}\right\rfloor
\]
tienen huecos
\[
p_{n+1}-p_n\in\{21,22\}.
\]
Los huecos cortos \(21\) retornan con separaciones
\[
6\quad\text{o}\quad7.
\]
Así aparece la jerarquía:
\[
\boxed{21/22\quad\longrightarrow\quad6/7.}
\]

Las ventanas clave son:
\[
22\epsilon=1.0066647923\ldots\Rightarrow Q_{22}\in\{1,2\},
\]
\[
132\epsilon=6.0399887540\ldots\Rightarrow Q_{132}\in\{6,7\},
\]
\[
1000\epsilon=45.7574905606\ldots\Rightarrow Q_{1000}\in\{45,46\}.
\]
Por tanto:
\[
1000-Q_{1000}\in\{955,954\}=\{900+55,900+54\}.
\]
Aparecen:
\[
45=S_{+,1},\qquad 54=\Delta_{\rm APP},\qquad55=\Delta_{\rm APP}+1.
\]

La carga angular desnuda es
\[
\alpha_0=\frac{\epsilon}{2\pi}.
\]
La renormalización HMT observada es
\[
\epsilon
=
2\pi\alpha
-\frac74\alpha^2
+\frac{\alpha^3}{2\pi}
+\frac{\alpha^4}{20}
-\frac{2}{21}\alpha^5
-\frac{\alpha^6}{46}
+\cdots.
\]
Con \(\alpha=\alpha_{\rm HMT}\), el residuo de seis términos cae a
\[
9.00\times10^{-18}.
\]
Al invertir la ecuación truncada de seis términos:
\[
\alpha^{-1}=137.0359990840000270\ldots,
\]
a \(2.70\times10^{-14}\) del inverso HMT.

Lectura:
\[
\boxed{e^-=\text{sección material estable asociada a la vacancia }z_t.}
\]
\[
\boxed{\text{corriente}=\text{deriva de fase de la red de vacancias}.}
\]
